% !TEX root = ../Monografia.tex
% ---------------------------------------------------------------------------------------------------------------------------------------------------------
% CAPÍTULO - RESULTADOS E DISCUSSÃO
% ---------------------------------------------------------------------------------------------------------------------------------------------------------

\chapter{Resultados e Discussão}

Este capítulo apresenta o estado final do pacote \texttt{collector\_flutter} e os resultados observados nos cenários de validação. A discussão está organizada em três frentes: módulos implementados, comportamento do pacote durante a execução do aplicativo de exemplo e limitações identificadas a partir das evidências coletadas.

% ---------------------------------------------------------------------------------------------------------------------------------------------------------
\section{Desenvolvimento do Pacote}
% ---------------------------------------------------------------------------------------------------------------------------------------------------------

O principal resultado é um pacote Dart/Flutter funcional, disponibilizado como código aberto e organizado em três camadas. Ao final do desenvolvimento, a solução reunia coleta de métricas, análise heurística, recomendações, painel integrado, exportação de dados e suíte automatizada de testes.

\subsection{Módulos Implementados}

A Tabela~\ref{tab:modulos-implementados} resume os módulos implementados e as evidências utilizadas para verificar sua entrega ao final do trabalho.

\begin{table}[H]
\centering
\caption{Módulos implementados e evidências de validação do \texttt{collector\_flutter}.}
\label{tab:modulos-implementados}
\small
\renewcommand{\arraystretch}{1.3}
\begin{tabular}{|p{4cm}|p{6cm}|p{4cm}|}
\hline
\textbf{Módulo} & \textbf{Funcionalidade} & \textbf{Evidência de validação} \\
\hline
\texttt{FrameDataSource}  & Coleta de FPS e detecção de \textit{jank} via \texttt{SchedulerBinding} & Cenário de \textit{rebuilds} intensivos \\
\hline
\texttt{MemoryDataSource} & Monitoramento de uso de memória RSS e \textit{heap} via \texttt{vm\_service} & Cenário de alocação intensiva \\
\hline
\texttt{HttpClientWrapper} & Interceptação de requisições HTTP e registro de eventos de rede & Cenário de chamadas simultâneas \\
\hline
\texttt{EventDataSource}  & Registro de eventos customizados definidos pelo desenvolvedor & Inspeção dos eventos exportados \\
\hline
\texttt{CpuDataSource}    & Coleta de uso de CPU por processo via \textit{platform channel} Android (\texttt{/proc/self/stat}) & Testes de canal e execução Android \\
\hline
\texttt{BatteryDataSource} & Monitoramento de bateria e estado de carregamento via \textit{platform channel} Android (\texttt{BatteryManager}) & Testes de canal e painel visual \\
\hline
\texttt{Analyzer}         & Análise heurística de métricas e geração de \texttt{AnalysisResult} & Casos unitários e cenários controlados \\
\hline
\texttt{Recommender}      & Geração de recomendações de otimização com níveis de severidade & Recomendações emitidas nos cenários \\
\hline
\texttt{ExportService}    & Exportação de sessão de telemetria em JSON e CSV & Testes de serialização \\
\hline
\texttt{CollectorBloc}    & Orquestração de estados e ciclo de coleta periódico & Atualização contínua do painel \\
\hline
\texttt{DashboardPage}    & Painel visual com gráficos em tempo real e lista de recomendações & Figuras~\ref{fig:evidencias-visuais-resultados} e Apêndice~\ref{ap:evidencias} \\
\hline
\texttt{ResourceCollector} & Ponto de entrada da biblioteca; gerencia ciclo de vida dos módulos & Execução do aplicativo de exemplo \\
\hline
Testes automatizados      & Suíte composta por 60 casos de teste distribuídos em quatro arquivos & Inspeção da suíte em \texttt{test/} e Apêndice~\ref{ap:testes} \\
\hline
\end{tabular}
\end{table}

\subsection{Estrutura da Biblioteca}

O sistema foi organizado em um pacote Dart publicável, com a seguinte estrutura de camadas:

\begin{itemize}
  \item \textbf{Camada de Dados (\textit{Data}):} contém as seis fontes de dados (\textit{data sources}) responsáveis pela coleta de métricas de quadros, memória, rede, eventos customizados, CPU e bateria, além do modelo de dados \texttt{TelemetryModel} e a implementação do repositório \texttt{TelemetryRepositoryImpl};

  \item \textbf{Camada de Domínio (\textit{Domain}):} define a interface \texttt{IResourceRepository}, a entidade de negócio \texttt{Telemetry} e os casos de uso responsáveis por coleta, análise e geração de recomendações, garantindo independência da camada de dados;

  \item \textbf{Camada de Apresentação (\textit{Presentation}):} implementa o gerenciamento de estados com \texttt{CollectorBloc}, o painel visual \texttt{DashboardPage}, com visualizações em tempo real e cartões de CPU e bateria, e a classe \texttt{ResourceCollector}, que funciona como fachada pública da biblioteca.
\end{itemize}

% ---------------------------------------------------------------------------------------------------------------------------------------------------------
\section{Validação da Arquitetura}
% ---------------------------------------------------------------------------------------------------------------------------------------------------------

Para verificar a aderência entre a arquitetura proposta e a implementação, foram realizadas inspeções de código e execuções manuais do pacote no aplicativo de exemplo (\texttt{example/}), que oferece botões de simulação de carga para os quatro cenários definidos na metodologia.

\subsection{Conformidade Arquitetural}

A implementação atende aos principais princípios arquiteturais adotados:

\begin{itemize}
  \item \textbf{Independência de camadas:} alterações nos \textit{data sources} permanecem restritas à camada de dados, preservando a camada de domínio;
  \item \textbf{Baixo acoplamento:} as dependências entre módulos são injetadas via construtores, sem uso de singletons globais (exceto o \texttt{HttpClientWrapper}, por necessidade de interceptação de rede);
  \item \textbf{Testabilidade:} a estrutura de casos de uso e repositórios permite substituição de dependências por \textit{mocks} em testes unitários, exercitada pela suíte composta por 60 casos de teste automatizados;
  \item \textbf{Extensibilidade:} novos tipos de métricas podem ser adicionados implementando um novo \textit{data source} sem modificar as camadas superiores, como demonstrado pela inclusão dos módulos de CPU e bateria.
\end{itemize}

\subsection{Comportamento Observado em Execução}

Durante as execuções manuais do aplicativo de exemplo em ambiente Android, incluindo a execução documentada no emulador \texttt{sdk gphone64 x86\_64} com Android 16 (API 36), foram registradas as seguintes observações:

\begin{itemize}
  \item O módulo de quadros registrou queda de FPS, aumento de P95 e frames longos quando o cenário de \textit{rebuilds} foi acionado;
  \item O cenário de memória gerou aumento mensurável de RSS nas leituras ADB, embora o painel tenha exibido 0{,}0 MB no snapshot capturado, limitação discutida na Seção~\ref{sec:cenario-memoria};
  \item A interceptação de rede registrou as requisições HTTP geradas pelo cenário de carga de rede, incluindo URL, método, código de status e latência;
  \item O módulo de CPU exibiu aumento de uso durante a carga combinada, chegando a 75{,}3\% no log registrado;
  \item As recomendações emitidas acompanharam os sintomas induzidos, principalmente nos cenários de rede e carga combinada.
\end{itemize}

% ---------------------------------------------------------------------------------------------------------------------------------------------------------
\section{Análise dos Cenários Experimentais}
% ---------------------------------------------------------------------------------------------------------------------------------------------------------

Os quatro cenários experimentais definidos na metodologia foram executados com o objetivo de verificar a operacionalidade do sistema e identificar inconsistências nas heurísticas. Os resultados apresentados a seguir são de natureza qualitativa.

A Tabela~\ref{tab:resultados-cenarios} sintetiza os principais sinais observados, as métricas acompanhadas e as evidências geradas durante a execução dos cenários. A Figura~\ref{fig:evidencias-visuais-resultados} mostra a diferença visual entre uma condição estável e uma condição com gargalos.
Na tela estável, os indicadores permanecem em faixas aceitáveis e a recomendação exibida é apenas informativa. Na tela com anomalias, o painel passa a destacar sintomas como aumento de eventos de rede, queda de fluidez e alertas de severidade, permitindo que o desenvolvedor relacione o comportamento observado ao cenário executado.

\begin{table}[H]
\centering
\caption{Síntese das evidências observadas nos cenários experimentais.}
\label{tab:resultados-cenarios}
\small
\renewcommand{\arraystretch}{1.3}
\footnotesize
\begin{tabular}{|p{2.8cm}|p{4cm}|p{3.8cm}|p{2.8cm}|}
\hline
\textbf{Cenário} & \textbf{Sinais observados} & \textbf{Métricas acompanhadas} & \textbf{Evidência gerada} \\
\hline
Reconstruções intensivas & 30 \textit{rebuilds}, FPS 55{,}4, P95 20{,}4 ms e 2 frames longos & FPS, \textit{frame time} e \textit{jank} & Registro do cenário e captura do painel \\
\hline
Chamadas de rede & 9 requisições HTTP, FPS 42{,}9, P95 139{,}5 ms e 4 frames longos & URL, método, código de status e duração & Eventos de rede no painel e exportação \\
\hline
Alocação de memória & RSS total de 236.480 KB e PSS total de 222.232 KB por ADB & RSS, \textit{heap} e tendência de memória & Captura do painel e leitura ADB \\
\hline
Carga combinada & 17 requisições, 54 \textit{rebuilds}, FPS 44{,}6, P95 37{,}7 ms e CPU 75{,}3\% & FPS, memória, rede, CPU e recomendações & Recomendações de frames e CPU \\
\hline
\end{tabular}
\end{table}

\begin{figure}[H]
  \centering
  \begin{minipage}{0.48\textwidth}
    \centering
    \includegraphics[height=0.50\textheight,keepaspectratio]{prototipoquandoestaok.jpeg}
  \end{minipage}
  \hfill
  \begin{minipage}{0.48\textwidth}
    \centering
    \includegraphics[height=0.50\textheight,keepaspectratio]{prototipoquandonaoestaok.jpeg}
  \end{minipage}
  \caption{Evidências visuais do painel durante a validação: à esquerda, condição estável; à direita, condição com gargalos detectados, eventos de rede e alertas ativos.}
  \label{fig:evidencias-visuais-resultados}
\end{figure}

\subsection{Cenário 1: Reconstruções Intensivas de Interface}

A simulação de \textit{rebuilds} frequentes aumentou o contador de reconstruções e alterou os indicadores de renderização. Na execução registrada no emulador, o log indicou 30 \textit{rebuilds}, FPS estimado de 55{,}4, P95 de 20{,}4 ms e 2 frames longos. A captura correspondente está apresentada na Figura~\ref{fig:ap-evidencia-rebuilds}, no Apêndice~\ref{ap:evidencias}.

\subsection{Cenário 2: Chamadas de Rede Simultâneas}

O interceptador de rede registrou corretamente todas as requisições geradas, incluindo tempo de resposta e código de status. Na execução registrada, o painel acumulou 9 requisições HTTP com respostas \texttt{200}; o log de evidência apontou FPS estimado de 42{,}9, P95 de 139{,}5 ms, 4 frames longos e recomendação associada à performance de frames. A captura correspondente está na Figura~\ref{fig:ap-evidencia-rede}. Observou-se que o buffer de eventos de rede, limitado a 1.000 entradas, foi suficiente para os cenários de teste executados.

\subsection{Cenário 3: Alocação Intensiva de Memória}
\label{sec:cenario-memoria}

O consumo de memória foi observado por meio do cenário de alocação controlada e de leitura complementar por ADB. Durante a alocação, o comando \texttt{adb shell dumpsys meminfo com.example.example} registrou RSS total de 236.480 KB e PSS total de 222.232 KB. Esse valor confirma crescimento real de memória no processo. No snapshot visual capturado, entretanto, o cartão de memória do painel apareceu como 0{,}0 MB, sinalizando uma limitação da leitura interna no emulador usado para a evidência. A captura correspondente está na Figura~\ref{fig:ap-evidencia-memoria}.

\subsection{Cenário 4: Carga Combinada}

A execução simultânea dos tipos de carga permitiu verificar que o painel continuou operável mesmo com múltiplos sinais sendo atualizados ao mesmo tempo. No cenário combinado, o log registrou 17 requisições, 54 \textit{rebuilds}, FPS estimado de 44{,}6, P95 de 37{,}7 ms e CPU em 75{,}3\%, resultando em recomendações de performance de frames e uso moderado de CPU. A captura correspondente está na Figura~\ref{fig:ap-evidencia-combinado}.

\subsection{Comparação Qualitativa com ADB}

As leituras por ADB foram utilizadas como apoio à interpretação dos cenários, não como cálculo formal de precisão. Em particular, o comando \texttt{dumpsys meminfo} forneceu leituras de memória do processo, enquanto \texttt{dumpsys gfxinfo} confirmou a presença da superfície Flutter, mas retornou contadores HWUI de frames zerados nessa execução. Por esse motivo, não foi calculado erro percentual de FPS a partir de \texttt{gfxinfo}.

\begin{table}[H]
\centering
\caption{Comparação qualitativa entre sinais do painel e leituras ADB.}
\label{tab:comparacao-adb}
\small
\renewcommand{\arraystretch}{1.25}
\begin{tabular}{|p{3cm}|p{5.4cm}|p{5.4cm}|}
\hline
\textbf{Cenário} & \textbf{Sinais do painel/log} & \textbf{Leitura ADB observada} \\
\hline
Reconstruções intensivas & 30 \textit{rebuilds}, FPS 55{,}4, P95 20{,}4 ms, CPU 22{,}1\%. & \texttt{meminfo}: RSS 199.208 KB; PSS 113.674 KB. \texttt{gfxinfo}: contadores HWUI indisponíveis para comparação. \\
\hline
Chamadas de rede & 9 requisições, FPS 42{,}9, P95 139{,}5 ms, 4 frames longos. & \texttt{meminfo}: RSS 213.380 KB; PSS 126.599 KB. \texttt{gfxinfo}: superfície ativa, sem série de frames comparável. \\
\hline
Alocação de memória & Cenário acionado no painel; leitura interna de memória indisponível no snapshot capturado. & \texttt{meminfo}: RSS 236.480 KB; PSS 222.232 KB durante a alocação. \\
\hline
Carga combinada & 17 requisições, 54 \textit{rebuilds}, FPS 44{,}6, P95 37{,}7 ms, CPU 75{,}3\%. & \texttt{gfxinfo}: superfície ativa, sem contadores HWUI utilizáveis; leitura posterior de \texttt{meminfo}: RSS 128.988 KB; PSS 107.293 KB. \\
\hline
\end{tabular}
\end{table}

\subsection{Evidências Complementares da Validação}

Além das observações por cenário, os resultados foram organizados em uma matriz de evidências para explicitar a relação entre o que foi testado, o artefato observado e a conclusão obtida. Essa organização liga os critérios definidos na metodologia aos resultados apresentados neste capítulo.

\begin{table}[H]
\centering
\caption{Matriz de evidências da validação funcional.}
\label{tab:matriz-evidencias}
\small
\renewcommand{\arraystretch}{1.3}
\begin{tabular}{|p{3.6cm}|p{5.4cm}|p{5.2cm}|}
\hline
\textbf{Dimensão avaliada} & \textbf{Evidência observada} & \textbf{Resultado obtido} \\
\hline
Coerência das métricas & Variação de FPS, \textit{frame time}, memória, rede e CPU acompanhando os cenários induzidos & Os sinais coletados acompanharam o tipo de carga aplicada em cada cenário \\
\hline
Pertinência das recomendações & Alertas emitidos pelo \texttt{Recommender} durante \textit{rebuilds}, chamadas de rede e carga combinada & As recomendações foram compatíveis com os sintomas provocados nos testes \\
\hline
Estabilidade do painel & Atualização contínua das métricas e mudança visual entre estado estável e estado de alerta & O painel permaneceu operável durante a coleta e destacou severidades distintas \\
\hline
Exportação e rastreabilidade & Sessões exportáveis em JSON e \textit{timings} de quadros em CSV pelo \texttt{ExportService} & Os dados coletados puderam ser inspecionados fora do painel \\
\hline
Confiabilidade interna & Suíte automatizada distribuída em quatro arquivos de teste & Cálculos, serialização, recomendações e \textit{fallbacks} dos canais nativos foram exercitados por testes automatizados \\
\hline
\end{tabular}
\end{table}

A Tabela~\ref{tab:distribuicao-testes} apresenta a distribuição dos 60 casos de teste automatizados por arquivo. Essa informação complementa a validação manual, pois mostra que os componentes centrais de análise, recomendação, exportação e coleta de CPU/bateria foram verificados também por testes unitários. O detalhamento nominal dos casos de teste está no Apêndice~\ref{ap:testes}.

\begin{table}[H]
\centering
\caption{Distribuição dos 60 casos de teste automatizados por arquivo.}
\label{tab:distribuicao-testes}
\small
\renewcommand{\arraystretch}{1.25}
\begin{tabular}{|p{4.4cm}|p{7.2cm}|p{1.8cm}|}
\hline
\textbf{Arquivo de teste} & \textbf{Representação gráfica} & \textbf{Casos} \\
\hline
Análise & \rule{6.0cm}{0.45em} & 24 \\
\hline
CPU e bateria & \rule{4.25cm}{0.45em} & 17 \\
\hline
Recomendações & \rule{3.0cm}{0.45em} & 12 \\
\hline
Exportação & \rule{1.75cm}{0.45em} & 7 \\
\hline
\textbf{Total} & \rule{6.0cm}{0.45em} & \textbf{60} \\
\hline
\end{tabular}
\end{table}

% ---------------------------------------------------------------------------------------------------------------------------------------------------------
\section{Evolução do Pacote: Limitações Superadas}
% ---------------------------------------------------------------------------------------------------------------------------------------------------------

As execuções da primeira fase do trabalho evidenciaram um conjunto de limitações técnicas, todas endereçadas ao longo do desenvolvimento:

\begin{enumerate}
  \item \textbf{Coleta de métricas de CPU:} implementada via \textit{platform channel} Android, no módulo \texttt{CpuDataSource}, que lê o arquivo \texttt{/proc/self/stat} e calcula a variação de tempo de CPU do processo entre amostras sucessivas. Em iOS, o canal retorna \texttt{-1.0} e o painel exibe ``N/D'', evitando falhas quando a métrica não está disponível;

  \item \textbf{Estimativa de consumo de bateria:} implementada via \textit{platform channel} Android, no módulo \texttt{BatteryDataSource}, que consulta o \texttt{BatteryManager} do sistema para obter o nível de carga (0--100\%) e o estado de carregamento. Um alerta de severidade \textit{medium} é gerado quando a bateria está abaixo de 20\% e o dispositivo não está carregando. Em iOS, o nível é obtido por \texttt{UIDevice.batteryLevel};

  \item \textbf{Persistência e exportação de dados:} o \texttt{ExportService} oferece exportação completa da sessão de telemetria em JSON estruturado, incluindo frames, memória, rede, CPU, bateria, análise e recomendações, além da exportação de timings de frames em CSV. Os dados são copiados para a área de transferência do sistema, permitindo integração com ferramentas externas de análise;

  \item \textbf{Testes automatizados:} a suíte é composta por 60 casos de teste distribuídos em quatro arquivos: \texttt{analyzer\_test.dart}, \texttt{recommender\_test.dart}, \texttt{export\_service\_test.dart} e \texttt{cpu\_battery\_test.dart}. Os testes verificam cálculos internos, detecção de anomalias, geração de recomendações, serialização e comportamento de \textit{fallback} dos \textit{platform channels};

  \item \textbf{Imprecisão no gráfico de frames:} a implementação inicial utilizava o instante atual como eixo X, em vez dos índices sequenciais dos quadros coletados. Esse defeito foi corrigido durante o desenvolvimento.
\end{enumerate}

% ---------------------------------------------------------------------------------------------------------------------------------------------------------
\section{Resumo do Capítulo}
% ---------------------------------------------------------------------------------------------------------------------------------------------------------

Este capítulo apresentou os resultados do desenvolvimento do \texttt{collector\_flutter}. A arquitetura de três camadas sustentou a inclusão dos módulos de CPU, bateria, exportação e testes sem alterar a organização central do pacote. Nos quatro cenários experimentais, o sistema permaneceu operável, registrou sinais compatíveis com as cargas induzidas e emitiu recomendações coerentes com os sintomas observados. As limitações encontradas, especialmente na leitura interna de memória e na ausência de série comparável via \texttt{gfxinfo}, delimitam o alcance da validação realizada.

\clearpage
